Pluripolar exceptions can also be removed within the logarithmic growth class. A pole barrier with the same growth allowance gives a direct repair of each competitor.

\begin{lemma}\label{lem:affine-logarithmic-pole-barrier}
For every pluripolar $E\subseteq\mathbb C^n$ and every $\delta>0$, there is a genuine $b\in\mathcal L_\delta$ with $b=-\infty$ on $E$. Given a non-pluripolar compact $K$, it can be chosen with $b\leq0$ on $K$.
\end{lemma}
\begin{proof}
Put $E_j=E\cap\overline B(0,j)$. Each $E_j$ is pluripolar on $\mathbb P^n$: it has local witnesses at every affine point, while a neighborhood of the hyperplane at infinity avoids this bounded set. By the compact global Josefson theorem, \cite[Theorem~1.1.5]{SGPT}, there is a genuine quasi-plurisubharmonic pole function for $E_j$. Rescaling its curvature bound and subtracting its supremum give
\[
 p_j\in\PSH(\mathbb P^n,\theta_1),\qquad
 p_j\leq0,\qquad p_j|_{E_j}=-\infty.
\]
The $L^1$ norms of these normalized potentials are uniformly bounded by \cref{thm:normalized-compactness}. The sums
\[
 p_N=\sum_{j=1}^N2^{-j}p_j
\]
decrease, have uniformly bounded $L^1$ norms, and remain $\theta_1$-plurisubharmonic because $\theta_1\geq0$ and the coefficients sum to at most one. \Cref{cor:decreasing-limit} gives a genuine canonical limit $p$. Every point of $E$ is a pole of one summand, hence $p=-\infty$ there. By \cref{lem:affine-projective-growth},
\[
 b=\delta(p+\chi/2)\in\mathcal L_\delta
\]
is the required affine pole function. Its supremum on $K$ is finite: compactness bounds it above, and non-pluripolarity excludes $b=-\infty$ on all of $K$. Subtracting this supremum gives $b\leq0$ on $K$.
\end{proof}
\index{Pluripolar set!logarithmic pole barrier}

\begin{corollary}\label{cor:affine-polar-obstacle-invariance}
The envelope $F_\delta$ is unchanged if $u|_K\leq\psi$ is imposed only quasi-everywhere, or only on $K\setminus E$ for any pluripolar $E\subseteq K$.
\end{corollary}
\begin{proof}
Let $u\in\mathcal L_\delta$ satisfy $u\leq\psi$ on $K\setminus E$, where $E$ is pluripolar. Choose $b$ from the lemma with this pole set and $b\leq0$ on $K$, and put $m=\inf_K\psi$. For $\epsilon>0$ define
\[
 u_\epsilon=\frac{u+\epsilon b+\epsilon m}{1+\epsilon}.
\]
This is a genuine member of $\mathcal L_\delta$. On $K\setminus E$, the inequalities $u\leq\psi$, $b\leq0$ and $m\leq\psi$ give $u_\epsilon\leq\psi$; on $E$ the pole term enforces the constraint. Thus $u_\epsilon\leq F_\delta$ and $u_\epsilon\to u$ almost everywhere. Canonical order gives $u\leq F_\delta$ everywhere. Taking the supremum and regularization proves the asserted equalities.

Equivalently, apply \cref{prop:compact-obstacle-polar-invariance} on $\mathbb P^n$ with the big background $\theta_\delta$ and the bounded obstacle $\psi-\delta\chi/2$, then use \eqref{eq:affine-projective-envelope}.
\end{proof}

These equalities concern regularizations. The raw envelope $V_\delta$ continues to use the obstacle at every point of $K$.
